% fig_fo1_2_crate_tail_family — peak 모양의 동역학 지연 L_V 가족 (Ch1 §6·§8·§9)
% 좌표 = eq:peakshape=(ξ_eq-ξ_lag)/L_V 의 인과 기억 적분(eq:lag)을 수치 평가(방전, w=1, U=0).
% L_V/w = 0(평형 종)·0.5·1.5·3.0. L_V∝|I| 라 C-rate 진행. 대상 라이브러리(ch1).
\begin{figure}[t]
\centering
\begin{tikzpicture}[x=0.46cm,y=12.4cm,>=Latex,font=\scriptsize]
% ---- 축 ----
\draw[->] (-6.4,0) -- (10.6,0) node[right] {$(V-U_j^{\,d})/w$};
\draw[->] (-6.4,0) -- (-6.4,0.275) node[above] {peak $(\xi_\eq-\xi_\mathrm{lag})/L_V$ [1/$w$]};
\foreach \xx in {-4,-2,0,2,4,6,8,10} {\draw (\xx,0.004)--(\xx,-0.004) node[below,font=\tiny] {\xx};}
\foreach \yy in {0.05,0.10,0.15,0.20,0.25} {\draw (-6.4,\yy)--(-6.7,\yy) node[left,font=\tiny] {\yy};}
% ---- L_V/w = 0 : 평형 종 (대칭, 점선 굵게) ----
\draw[densely dotted,very thick] plot[smooth] coordinates {(-6.0000,0.0025) (-5.3333,0.0048) (-4.6667,0.0092) (-4.0000,0.0177) (-3.3333,0.0333) (-2.6667,0.0607) (-2.0000,0.1050) (-1.3333,0.1651) (-0.6667,0.2242) (0.0000,0.2500) (0.6667,0.2242) (1.3333,0.1651) (2.0000,0.1050) (2.6667,0.0607) (3.3333,0.0333) (4.0000,0.0177) (4.6667,0.0092) (5.3333,0.0048) (6.0000,0.0025)};
% ---- L_V/w = 0.5 (얇은 실선) ----
\draw[thin] plot[smooth] coordinates {(-6.0000,0.0016) (-5.3333,0.0032) (-4.6667,0.0062) (-4.0000,0.0119) (-3.3333,0.0226) (-2.6667,0.0419) (-2.0000,0.0745) (-1.3333,0.1232) (-0.6667,0.1815) (0.0000,0.2274) (0.3333,0.2371) (0.6667,0.2353) (1.3333,0.2016) (2.0000,0.1471) (2.6667,0.0952) (3.3333,0.0566) (4.0000,0.0319) (4.6667,0.0173) (5.3333,0.0092) (6.0000,0.0048) (7.0000,0.0018) (8.0000,0.0007) (9.0000,0.0002) (10.0000,0.0001)};
% ---- L_V/w = 1.5 (파선) ----
\draw[thick,densely dashed] plot[smooth] coordinates {(-6.0000,0.0010) (-5.3333,0.0019) (-4.6667,0.0037) (-4.0000,0.0072) (-3.3333,0.0137) (-2.6667,0.0255) (-2.0000,0.0462) (-1.3333,0.0785) (-0.6667,0.1214) (0.0000,0.1647) (0.6667,0.1918) (1.0000,0.1955) (1.6667,0.1838) (2.3333,0.1550) (3.0000,0.1210) (3.6667,0.0894) (4.3333,0.0636) (5.0000,0.0441) (5.6667,0.0300) (6.3333,0.0201) (7.0000,0.0133) (8.0000,0.0071) (9.0000,0.0037) (10.0000,0.0020)};
% ---- L_V/w = 3.0 (굵은 실선) ----
\draw[very thick] plot[smooth] coordinates {(-6.0000,0.0006) (-5.3333,0.0012) (-4.6667,0.0023) (-4.0000,0.0045) (-3.3333,0.0086) (-2.6667,0.0161) (-2.0000,0.0293) (-1.3333,0.0504) (-0.6667,0.0796) (0.0000,0.1119) (0.6667,0.1375) (1.3333,0.1490) (2.0000,0.1458) (2.6667,0.1328) (3.3333,0.1153) (4.0000,0.0972) (4.6667,0.0804) (5.3333,0.0657) (6.0000,0.0533) (6.6667,0.0430) (7.3333,0.0346) (8.0000,0.0278) (8.6667,0.0223) (9.3333,0.0179) (10.0000,0.0144)};
% ---- 정점 표시 ----
\fill (0,0.25) circle (1.3pt);
\fill (0.333,0.2371) circle (1.1pt);
\fill (1.0,0.1955) circle (1.1pt);
\fill (1.333,0.1490) circle (1.1pt);
% ---- 곡선 라벨 ----
\node[anchor=west,font=\tiny] at (-6.2,0.263) {$L_V/w{=}0$ (평형 종, $|I|\!\to\!0$)};
\node[anchor=south west,font=\tiny] at (0.5,0.239) {$0.5$};
\node[anchor=south,font=\tiny] at (1.2,0.199) {$1.5$};
\node[anchor=south west,font=\tiny] at (1.6,0.150) {$3.0$};
% ---- 진행 화살표 ----
\draw[->,black!55] (-1.8,0.235) to[bend left=12] (2.3,0.135);
\node[anchor=west,align=left,font=\tiny] at (2.7,0.205)
  {$L_V\!\uparrow$ ($|I|\!\uparrow$):\\정점 낮아지고\\꼬리 높은 $V$ 로 늘어남};
\draw[->,gray] (6.5,0.055) -- (9.2,0.030);
\node[anchor=west,font=\tiny,gray] at (6.2,0.078) {인과 꼬리 $\to$ 높은 $V$};
\end{tikzpicture}
\caption{동역학 지연 길이 $L_{V,j}$ 가 평형 종을 꼬리로 펴는 모습 --- 좌표는 peak 모양
$(\xi_{\eq,j}-\xi_{\mathrm{lag},j})/L_{V,j}$(식~\eqref{eq:peakshape}, 방전)의 인과 기억
적분(식~\eqref{eq:lag})을 무차원 폭 $w$ 로 규격화해($V,U_j^{\,d}$ 를 $w$ 단위로) 식 그대로 수치 평가한
것이다. $L_V/w{=}0$ 점선이 $|I|\!\to\!0$ 평형 종(정점 $1/(4w){=}0.25$, 대칭, 식~\eqref{eq:tail-limit} 의
극한)이고, $L_V\!\propto\!|I|$(식~\eqref{eq:LV})가 커질수록 정점이 $0.237\!\to\!0.196\!\to\!0.149$ 로 낮아지며
높은 $V$ 쪽으로 치우친 인과 꼬리가 자란다 --- 서론이 든 ``저율$\cdot$고율에서 꼬리가 길어져 이웃 peak 와
겹친다''는 관측의 C-rate 갈래가 이 한 식의 $L_V$ 의존이다. 면적(용량 $Q_j$)은 보존된다.}
\label{fig:cand-fo1-2}
\end{figure}
